\begin{figure}[H]
\centering
\caption*{\textbf{Figura 8 - Sequência de uma execução completa da MOSAIC 0.2}}
\begin{tikzpicture}[node distance=4.5mm]
\node[mosaicbox, text width=12.5cm] (s1)
  {1. Usuário \(\rightarrow\) planner: pedido, contexto e restrições};
\node[mosaicbox, text width=12.5cm, below=of s1] (s2)
  {2. Planner \(\leftrightarrow\) catálogo: P0, hints body-aware e uma revisão P1};
\node[mosaicbox, text width=12.5cm, below=of s2] (s3)
  {3. Scheduler \(\rightarrow\) router: objetivo pronto, \texttt{doneWhen} e estado projetado};
\node[mosaicbox, text width=12.5cm, below=of s3] (s4)
  {4. Router \(\leftrightarrow\) catálogo: retrieve, rerank pelo body e bundle ordenado};
\node[mosaicbox, text width=12.5cm, below=of s4] (s5)
  {5. Bundle + registro de tools \(\rightarrow\) executor: menu composto e bodies delimitados};
\node[mosaicbox, text width=12.5cm, below=of s5] (s6)
  {6. Modelo \(\leftrightarrow\) tools: o runtime correlaciona cada retorno e preserva a ordem};
\node[mosaicbox, text width=12.5cm, below=of s6] (s7)
  {7. Modelo \(\rightarrow\) runtime: decisão semântica; o runtime anexa todas as observações
  e materializa o resultado do nó};
\node[mosaicbox, text width=12.5cm, below=of s7] (s8)
  {8. Se houver revisão, o planner recebe a suposição invalidada, o efeito solicitado
  e todo o conjunto de observações, sem IDs do provedor};
\node[mosaicbox, text width=12.5cm, below=of s8] (s9)
  {9. Quando os objetivos entregáveis terminam, o compositor empacota os resultados
  sem nova chamada ao modelo};
\foreach \a/\b in {s1/s2,s2/s3,s3/s4,s4/s5,s5/s6,s6/s7,s7/s8,s8/s9}
  \draw[mosaicarrow] (\a) -- (\b);
\end{tikzpicture}
\par\smallskip\raggedright\small Fonte: elaboração própria (2026).
\end{figure}
